\chapter{Introducción}
\chaptermark{Introducción}
\label{intro}

\section{Contexto General}

El sueño es un proceso biológico activo y altamente organizado, fundamental para la homeostasis neurofisiológica. Lejos de ser un estado uniforme de inactividad, el cerebro dormido transita cíclicamente por distintas fases (Vigilia, Sueño Ligero S1 y S2, Sueño de Ondas Lentas SWS, y Sueño de Movimientos Oculares Rápidos REM), cada una con funciones específicas de restauración metabólica y consolidación de memoria.

Históricamente, el estándar de oro para estudiar esta arquitectura ha sido la polisomnografía y la estadificación manual basada en las reglas de \cite{rechtschaffen1968manual}. Tradicionalmente, el análisis clínico y de investigación ha dependido de la \textbf{estadística descriptiva global}, resumiendo noches enteras de 8 horas en porcentajes totales (ej. ``20\% de REM'', ``50\% de S2''). Este enfoque asume implícitamente que la estructura del sueño es estática y que todas las épocas clasificadas bajo una misma etiqueta (como la Fase 2) son funcionalmente idénticas.

\section{Planteamiento del Problema}

El problema central que aborda esta tesis es la \textbf{insuficiencia de las métricas globales y estáticas} para capturar la dinámica secuencial del sueño.
La Fase 2 (S2), que constituye aproximadamente el 50\% del tiempo total de sueño en adultos, ha sido considerada históricamente como una etapa de ``transición pasiva'' o un simple puente entre estados. Sin embargo, la evidencia biológica sugiere que el sueño es un proceso no-estacionario y dependiente del contexto.

Para ilustrar esta limitación, la Figura \ref{fig:problema_conceptual} contrasta la visión clásica con la realidad dinámica del sueño.

\begin{figure}[htbp]
    \centering
    \begin{subfigure}[b]{0.62\textwidth}
        \centering
        \includegraphics[width=\textwidth]{intro_estatico_donut.png}
        \caption{Enfoque Estático: Porcentajes globales.}
        \label{fig:estatico}
    \end{subfigure}

    \vspace{0.6em}

    \begin{subfigure}[b]{0.92\textwidth}
        \centering
        \includegraphics[width=\textwidth]{intro_dinamico_hipno.png}
        \caption{Enfoque Dinámico: Trayectoria secuencial.}
        \label{fig:dinamico}
    \end{subfigure}
    \caption{El Paradigma Estático vs. Dinámico. (A) Las métricas globales (porcentajes) ocultan la estructura temporal. (B) El análisis secuencial revela la intencionalidad y trayectoria de las transiciones.}
    \label{fig:problema_conceptual}
\end{figure}

Al reducir la complejidad del hipnograma a promedios globales, se pierde información crítica sobre la \textbf{intencionalidad} de las transiciones. Actualmente, no existen herramientas estandarizadas que permitan diferenciar si una época de S2 es un estado de ``profundización'' (preparándose para entrar a SWS) o un estado de ``activación'' (preparándose para REM). Esta ceguera metodológica impide comprender alteraciones sutiles en la arquitectura del sueño que no afectan los porcentajes totales, pero sí la continuidad y funcionalidad del ciclo.

\section{Justificación de la Investigación}

Esta investigación se justifica por la necesidad de incorporar herramientas modernas de \textbf{Inteligencia Artificial Interpretativa} al análisis de señales biomédicas. Mientras que el estado del arte en Deep Learning se ha enfocado en la clasificación automática (``cajas negras''), esta tesis propone un cambio de paradigma hacia el análisis de la dinámica.

Nuestros análisis (ver Capítulo \ref{resultados}) revelaron que la duración de los bloques de Fase 2 no sigue una distribución normal, sino una \textbf{distribución de cola pesada}, en la que la media (12.10 épocas) y la mediana (5 épocas) divergen marcadamente. Esto implica que no existe una ``duración promedio'' representativa y que los métodos estadísticos tradicionales (basados en la media y desviación estándar) resultan inadecuados.
La aplicación de técnicas de \textbf{Procesamiento de Lenguaje Natural (NLP)}, específicamente modelos de Embeddings, permite superar esta limitación. Al tratar las fases del sueño como ``palabras'' y los ciclos como ``oraciones'', podemos interrogar la estructura gramatical latente del hipnograma, abriendo la puerta a nuevos biomarcadores digitales basados en la trayectoria y no solo en la cantidad de sueño.

Finalmente, la relevancia clínica de esta propuesta radica en su potencial como \textbf{biomarcador temprano}. Muchas patologías del sueño y neurodegenerativas alteran la micro-arquitectura del sueño antes de que los porcentajes globales (el hipnograma macroscópico) muestren cambios evidentes. Si nuestra hipótesis es correcta y la Fase 2 posee subtipos funcionales, la pérdida de esta diferenciación dinámica podría servir como una señal de alerta precoz de fragmentación del sueño, invisible para las métricas actuales.

\section{Preguntas de Investigación e Hipótesis}

\subsection*{Interrogantes de Investigación}

\begin{enumerate}
    \item ¿Existe información en la estructura secuencial de los hipnogramas que permita caracterizar la dinámica del sueño (S2, SWS, REM) en los distintos ciclos de sueño?
    \item ¿Pueden los hipnogramas tratarse como secuencias simbólicas que, al segmentarse en unidades de varias épocas, exhiben las regularidades estadísticas propias de un lenguaje natural?
    \item ¿Es posible anticipar su transición a otras fases (SWS o REM) a partir de los vectores de incrustación?
\end{enumerate}

\subsection*{Hipótesis}

\begin{itemize}
    \item \textbf{H1 (Estructural):} El sueño de un participante sano no es un proceso aleatorio; sigue una estructura particular no fragmentada.
    \item \textbf{H2 (Lingüística):} Los hipnogramas son secuencias simbólicas que pueden ser tratadas mediante principios homólogos a los de un lenguaje natural.
    \item \textbf{H3 (Dinámica):} La representación vectorial de la Fase 2 presenta una bifurcación topológica en el espacio latente que correlaciona con su destino (SWS vs. REM), permitiendo anticipar la transición antes de que ocurra clínicamente.
\end{itemize}

\section{Objetivos}

\textbf{Objetivo General:}
Modelar la dinámica de las transiciones de sueño mediante un enfoque híbrido de Cadenas de Markov y algoritmos de Embeddings (Word2Vec), para caracterizar la heterogeneidad funcional de la Fase 2.

\textbf{Objetivos Específicos:}
\begin{enumerate}
    \item Cuantificar la estabilidad y las probabilidades de transición explícitas entre fases utilizando Matrices de Markov, diferenciando por ciclos de sueño para respetar la no-estacionariedad.
    \item Determinar la escala temporal de tokenización a la que las secuencias del hipnograma reproducen las regularidades estadísticas de un lenguaje natural, mediante un criterio multi-objetivo basado en la ley de Zipf.
    \item Entrenar un modelo de representación vectorial (Word2Vec Skip-Gram) sobre secuencias de hipnogramas para capturar las dependencias contextuales de largo plazo.
    \item Implementar técnicas de reducción de dimensionalidad (UMAP) y métricas de similitud coseno para visualizar y medir la trayectoria dinámica de la Fase 2 en el espacio latente.
    \item Validar si existe una ``bifurcación'' en la trayectoria de la Fase 2 que anticipe la entrada a estados de sueño profundo o REM.
\end{enumerate}

\section{Alcance y Delimitaciones}

El núcleo metodológico de este estudio se construye sobre hipnogramas de adultos sanos, que definen una referencia normativa de la arquitectura del sueño a partir de la cual se modela su dinámica. Metodológicamente, la resolución temporal del análisis ($N=6$ épocas) se fija mediante un criterio multi-objetivo que combina el ajuste a la ley de Zipf, la calidad de dicho ajuste y la cobertura de los bloques de Fase 2, sin recurrir a la calibración por un ciclo particular.

\FloatBarrier
\section{Estructura de la Tesis}

El presente trabajo se organiza en cinco capítulos:
\begin{itemize}
    \item \textbf{Capítulo 1: Introducción.} Presenta el contexto, la problemática y los objetivos del estudio.
    \item \textbf{Capítulo 2: Marco Teórico.} Revisa la fisiología del sueño, las limitaciones de los modelos estocásticos clásicos y fundamenta el uso de algoritmos de NLP y Manifold Learning.
    \item \textbf{Capítulo 3: Metodología.} Describe el procesamiento de datos, la configuración del modelo Word2Vec y los parámetros de visualización con UMAP.
    \item \textbf{Capítulo 4: Resultados.} Expone los hallazgos sobre la estabilidad markoviana y presenta las proyecciones vectoriales que evidencian la dinámica de la Fase 2.
    \item \textbf{Capítulo 5: Discusión y Conclusiones.} Interpreta los resultados a la luz de la fisiología del sueño, discute las implicaciones clínicas y propone líneas de trabajo futuro.
\end{itemize}
